\section*{Capítulo \ref{ch:b3:adaptive-immunity} --- Imunidade adaptativa e vacinação}

\begin{solution}{exo:b3:adaptive-immunity:1}
Um \omterm{def:b3:adaptive-immunity:lymphocytes}{linfócito}, uma especificidade (Nossal e Lederberg: cada célula
secretava \omterm{def:b3:adaptive-immunity:antibody}{anticorpo} contra um só \omterm{def:b3:adaptive-immunity:lymphocytes}{antígeno}); o receptor é feito antes de
o \omterm{def:b3:adaptive-immunity:lymphocytes}{antígeno} ser encontrado (Tonegawa: o gene é rearranjado na \omterm{def:b3:adaptive-immunity:lymphocytes}{célula B},
e não em resposta ao \omterm{def:b3:adaptive-immunity:lymphocytes}{antígeno}); o \omterm{def:b3:adaptive-immunity:lymphocytes}{antígeno} seleciona e expande os clones
compatíveis em células efetoras e de memória; os clones autorreativos
são deletados durante o desenvolvimento.
\end{solution}

\begin{solution}{exo:b3:adaptive-immunity:2}
Duas cadeias pesadas e duas leves, cada uma com um domínio variável na
ponta e domínios constantes atrás; dois sítios de ligação idênticos
formados por seis alças hipervariáveis; uma haste Fc que outras células
e o \omterm{def:b3:innate-immunity:complement}{complemento} leem. IgM: primeiro \omterm{def:b3:adaptive-immunity:antibody}{anticorpo}, pentâmero, ativação do
\omterm{def:b3:innate-immunity:complement}{complemento}. IgG: principal \omterm{def:b3:adaptive-immunity:antibody}{anticorpo} do sangue e dos tecidos,
\omterm{def:b3:innate-immunity:complement}{opsonização}, neutralização, atravessa a placenta. IgA: dímero secretado
sobre as superfícies mucosas e no leite. IgE: presa aos mastócitos,
contra parasitas, mediadora da alergia. IgD: nas \omterm{def:b3:adaptive-immunity:lymphocytes}{células B} virgens, de
função conhecida menor.
\end{solution}

\begin{solution}{exo:b3:adaptive-immunity:3}
Classe~I: todas as células nucleadas; peptídeos de proteínas
citosólicas cortados pelo proteassomo; 8--10 resíduos; lidos pelas
\omterm{def:b3:adaptive-immunity:response}{células T citotóxicas} CD8. Classe~II: \omterm{def:b3:innate-immunity:innate}{células dendríticas}, \omterm{def:b3:innate-immunity:innate}{macrófagos},
\omterm{def:b3:adaptive-immunity:lymphocytes}{células B}; peptídeos de proteínas captadas e digeridas em \omterm{def:b3:membrane-traffic:endocytosis}{endossomos};
13--25 resíduos; lidos pelas \omterm{def:b3:adaptive-immunity:response}{células T auxiliares} CD4.
\end{solution}

\begin{solution}{exo:b3:adaptive-immunity:4}
$R_{0}$ é o número médio de casos que um caso produz numa população
inteiramente suscetível; o limiar é $p_{c} = 1 - 1/R_{0}$. Caxumba,
\qty{80}{\%}; coqueluche, \qty{92}{\%}; gripe de 1918, \qty{60}{\%}.
\end{solution}

\begin{solution}{exo:b3:adaptive-immunity:5}
$50\times 20\times 5 = 5000$ cadeias pesadas, $60\times 4 = 240$ leves:
$1.2\times 10^{6}$ combinações; com a \omterm{def:b3:adaptive-immunity:antibody}{diversidade juncional}, $1.2\times
10^{10}$ --- dez vezes o número de \omterm{def:b3:adaptive-immunity:lymphocytes}{células B}, de modo que o animal
expressa no máximo um décimo do que poderia fazer, e quase toda célula
carrega um receptor único.
\end{solution}

\begin{solution}{exo:b3:adaptive-immunity:6}
$1000\times 10^{-3} = 1$ mutação por célula por divisão; $20$ em vinte
divisões. Entre $10^{5}$ células, $2\times 10^{6}$ eventos de mutação
caem sobre $3000$ substituições únicas possíveis: cada uma é amostrada
centenas de vezes, e ainda as combinações duplas e triplas --- o \omterm{def:b3:adaptive-immunity:germinal-centre}{centro
germinativo} explora exaustivamente a vizinhança do receptor.
\end{solution}

\begin{solution}{exo:b3:adaptive-immunity:7}
As \omterm{def:b3:adaptive-immunity:lymphocytes}{células T} são selecionadas no \omterm{def:b3:adaptive-immunity:tolerance}{timo} para reconhecer peptídeos nos
alelos de MHC do próprio hospedeiro; os alelos de outra pessoa têm
fendas de formato diferente, que seguram peptídeos diferentes em
orientações diferentes, de modo que o mesmo peptídeo viral não é visto.
As moléculas de MHC estranhas de um enxerto são elas próprias
reconhecidas por grande parte das \omterm{def:b3:adaptive-immunity:lymphocytes}{células T} como ``MHC próprio com um
peptídeo esquisito'', daí a rejeição. O polimorfismo persiste porque uma
população com muitos alelos apresenta uma amostra mais ampla dos
peptídeos de qualquer patógeno e nenhum patógeno consegue escapar de
todos; os heterozigotos apresentam o dobro do repertório e são
favorecidos.
\end{solution}

\begin{solution}{exo:b3:adaptive-immunity:8}
(a) Sem \omterm{def:b3:adaptive-immunity:antibody}{recombinação V(D)J}: nenhuma \omterm{def:b3:adaptive-immunity:lymphocytes}{célula B} ou T, imunodeficiência
combinada grave. (b) Sem AIRE: as células tímicas não conseguem exibir
os \omterm{def:b3:adaptive-immunity:lymphocytes}{antígenos} teciduais, as \omterm{def:b3:adaptive-immunity:lymphocytes}{células T} autorreativas escapam,
\omterm{def:b3:adaptive-immunity:tolerance}{autoimunidade} contra muitos órgãos. (c) Sem \omterm{def:b3:adaptive-immunity:tolerance}{FoxP3}: nenhuma \omterm{def:b3:adaptive-immunity:response}{célula T
reguladora}, \omterm{def:b3:adaptive-immunity:tolerance}{autoimunidade} e alergia fatais de início precoce. (d) Sem
AID: nenhuma \omterm{def:b3:adaptive-immunity:antibody}{troca de classe} nem hipermutação --- IgM abundante, quase
nenhuma IgG ou IgA, \omterm{def:b3:adaptive-immunity:antibody}{anticorpos} de baixa afinidade, infecções de
repetição. (e) Sem \omterm{def:b3:adaptive-immunity:lymphocytes}{células T} CD4: nenhuma ajuda para as \omterm{def:b3:adaptive-immunity:lymphocytes}{células B} nem
para os \omterm{def:b3:innate-immunity:innate}{macrófagos}, respostas de \omterm{def:b3:adaptive-immunity:antibody}{anticorpo} pobres e suscetibilidade a
infecções oportunistas --- a imunodeficiência da AIDS.
\end{solution}

\begin{solution}{exo:b3:adaptive-immunity:9}
$p_{c} = 1 - 1/6 = 0.83$; cobertura $0.83/0.9 = 93\,\%$. Com cobertura
de \qty{80}{\%} a fração imune é $0.72$ e $R_{\text{eff}} = 6\times 0.28
= 1.7$. Entre os suscetíveis a epidemia se espalha como se $R_{0}$
fosse $1.7$ (os imunes apenas diluem os contatos), e a relação de
tamanho final $s_{\infty} - \ln s_{\infty}/1.7 = 1$ dá $s_{\infty}
\approx 0.31$: cerca de \qty{69}{\%} dos suscetíveis, \qty{19}{\%} da
população, são infectados.
\end{solution}

\begin{solution}{exo:b3:adaptive-immunity:10}
De $\mathrm{d}s/\mathrm{d}t = -\beta si$ e $\mathrm{d}i/\mathrm{d}t
= \beta si - \gamma i$ vem $\mathrm{d}i/\mathrm{d}s = -1 + \gamma/(\beta s)$,
de modo que $i + s - (\ln s)/R_{0}$ se conserva; com $s = 1$, $i = 0$ no
início e $i = 0$, $s = s_{\infty}$ no fim, $s_{\infty} - \ln
s_{\infty}/R_{0} = 1$. Soluções: $R_{0} = 1.5$, $s_{\infty} = 0.42$
(\qty{58}{\%} infectados); $2$, $0.20$ (\qty{80}{\%}); $3$, $0.06$
(\qty{94}{\%}); $5$, $0.007$ (\qty{99.3}{\%}). Não todo mundo: à medida
que os suscetíveis se esgotam, o número efetivo de reprodução cai abaixo
de um enquanto alguns ainda restam, e a epidemia se apaga antes de
alcançá-los.
\end{solution}

\begin{solution}{exo:b3:adaptive-immunity:11}
$1.5^{n} = 1000$: $n = \ln 1000/\ln 1.5 \approx 17$ melhorias
sucessivas. Com uma mutação por célula por divisão e uma em cem
melhorando, é preciso triar ao menos cem células por rodada para achar
uma; na prática, milhares, já que metade das mutações destrói a ligação
e a seleção é ruidosa. Dezessete rodadas de duas divisões cada levam
cerca de nove dias --- as duas a três semanas observadas, dadas as
ineficiências. Os ciclos separam a mutação (zona escura) do teste (zona
clara), de modo que cada variante é julgada contra o \omterm{def:b3:adaptive-immunity:lymphocytes}{antígeno} antes de
ter permissão para se multiplicar, exatamente como um algoritmo
genético alterna variação e seleção.
\end{solution}

\begin{solution}{exo:b3:adaptive-immunity:12}
Mulheres: vários genes imunes ficam no cromossomo X, alguns (TLR7)
escapam da inativação e são expressos das duas cópias, e os estrógenos
aumentam a sobrevivência das \omterm{def:b3:adaptive-immunity:lymphocytes}{células B} --- teste: homens com um X extra
(Klinefelter) deveriam ter risco de lúpus semelhante ao feminino, e têm.
\omterm{def:b3:adaptive-immunity:mhc}{HLA}: os alelos associados apresentam bem os peptídeos próprios
relevantes, ou deixam de apresentá-los no \omterm{def:b3:adaptive-immunity:tolerance}{timo}, de modo que os clones
não são deletados --- teste: camundongos \omterm{def:b3:genetic-engineering:transgenic}{transgênicos} para o alelo
humano desenvolvem a doença quando recebem o \omterm{def:b3:adaptive-immunity:lymphocytes}{antígeno}. A alta: menos
exposição microbiana precoce e \omterm{def:b3:microbiomes:symbiosis}{microbiomas} alterados (\omterm{def:b3:bacteriology:antibiotics}{antibióticos},
dieta, parto cesáreo) dão menos \omterm{def:b3:adaptive-immunity:response}{células T reguladoras} e mais
\omterm{def:b3:innate-immunity:inflammation}{inflamação} --- testes: camundongos propensos a diabetes criados livres
de germes ou sob \omterm{def:b3:bacteriology:antibiotics}{antibióticos} desenvolvem mais diabetes, e os filhos de
imigrantes adquirem a incidência do novo país em uma geração.
\end{solution}

\begin{solution}{pb:b3:adaptive-immunity:1}
\textbf{1.} $6000\times 320 = 1.9\times 10^{6}$; $\times 3\times 10^{4}
= 5.8\times 10^{10}$.
\textbf{2.} $10^{11}$ células contra $5.8\times 10^{10}$ possíveis: um
corpo poderia abrigar cada receptor cerca de duas vezes, mas os clones
são desiguais e a maioria das sequências simplesmente nunca é feita, de
modo que cada corpo expressa uma amostra grande mas incompleta --- e
duas pessoas compartilham apenas uma pequena fração de seus receptores.
\textbf{3.} $10^{11}/10^{5} = 10^{6}$ células; num \omterm{def:b3:adaptive-immunity:lymphocytes}{linfonodo} de
$10^{8}$, cerca de $1000$.
\textbf{4.} $10^{9}/10^{5} = 10^{4}$ precursores --- cem vezes menos; o
recém-nascido tem ainda folículos e \omterm{def:b3:innate-immunity:innate}{células dendríticas} imaturos,
nenhuma memória e ajuda limitada das \omterm{def:b3:adaptive-immunity:lymphocytes}{células T}, de modo que as respostas
são lentas e pequenas, e a IgG materna faz a ponte.
\textbf{5.} Ela abrange as junções V--D e D--J, onde nucleotídeos são
removidos e acrescentados ao acaso: sua sequência não está codificada
em lugar nenhum do genoma. Situada entre as duas cadeias, no centro do
sítio de ligação, é a alça mais frequentemente em contato com o epítopo.
\textbf{6.} Antes do \omterm{def:b3:adaptive-immunity:lymphocytes}{antígeno}, a diversidade precisa ser ampla para que
algo ligue o que quer que venha; depois do \omterm{def:b3:adaptive-immunity:lymphocytes}{antígeno}, o esforço rende
mais variando os poucos receptores que já se sabe que ligam. Uma busca
grosseira seguida de uma fina é muito mais eficiente que qualquer uma
das duas sozinha.
\textbf{7.} Sete dias são $21$ divisões: $100\times 2^{21} = 2.1\times
10^{8}$ células; \qty{30}{\%}, $6.3\times 10^{7}$ \omterm{def:b3:adaptive-immunity:response}{plasmócitos}.
\textbf{8.} $6.3\times 10^{7}\times 2000\times \num{86400} = 1.1\times
10^{16}$ moléculas por dia; $\times 1.5\times 10^{5}\times 1.66\times
10^{-24}$ g $= \qty{2.7}{mg}$; em \qty{3}{L}, \qty{0.9}{mg/L} por dia.
\textbf{9.} \qty{27}{mg} em \qty{3}{L}: \qty{9}{mg/L} $= \qty{0.009}{g/L}$,
um milésimo da IgG total --- o \omterm{def:b3:adaptive-immunity:antibody}{anticorpo} contra um único \omterm{def:b3:adaptive-immunity:lymphocytes}{antígeno} é uma
pequena fração do todo.
\textbf{10.} Cem vezes são $\log_{2}100 = 6.6$ meias-vidas, cerca de
$140$ dias. Anos de \omterm{def:b3:adaptive-immunity:antibody}{anticorpo} vêm dos \omterm{def:b3:adaptive-immunity:response}{plasmócitos} de vida longa da
medula, que secretam continuamente, e das \omterm{def:b3:adaptive-immunity:response}{células de memória}
reestimuladas na reexposição.
\textbf{11.} Nove divisões: $10^{5}\times 512 = 5\times 10^{7}$ células no
dia 3 --- mil vezes as $5\times 10^{4}$ da primária no dia 3, e um
quarto do que a primária só alcançou no dia 7.
\textbf{12.} O VDJ rearranjado fica ao lado de C$_{\mu}$, de modo que o
primeiro transcrito é IgM; trocar para C$_{\gamma}$ exige ajuda das
\omterm{def:b3:adaptive-immunity:lymphocytes}{células T} e a AID, o que leva dias. As \omterm{def:b3:adaptive-immunity:response}{células de memória} já trocaram,
de modo que sua progênie faz IgG de imediato.
\textbf{13.} $1000\times 10^{-3} = 1$ por divisão; $20$ ao longo de
vinte.
\textbf{14.} Uma em cem; $10^{4}\times 0.01 = 100$ células melhoradas
por divisão.
\textbf{15.} $1.5^{n} = 1000$: $n \approx 17$; $17\times \qty{12}{h}
\approx 8.5$ dias.
\textbf{16.} Metade das filhas perde a ligação: sem seleção, $0.5^{20}
\approx 10^{-6}$ da linhagem ainda ligaria depois de vinte divisões. A
zona clara precisa remover os perdedores a cada ciclo para que só
receptores funcionais e melhorados sigam adiante.
\textbf{17.} A primeira dose deixa \omterm{def:b3:adaptive-immunity:lymphocytes}{células B} de memória de afinidade
elevada e classe trocada; a segunda dose semeia com elas novos \omterm{def:b3:adaptive-immunity:germinal-centre}{centros
germinativos}, e a hipermutação e a seleção recomeçam de um ponto de
partida mais alto, de modo que os \omterm{def:b3:adaptive-immunity:antibody}{anticorpos} que emergem são melhores
ainda.
\textbf{18.} Algumas pessoas infectadas acabam fazendo \omterm{def:b3:adaptive-immunity:antibody}{anticorpos} contra
sítios conservados que o \omterm{def:b3:virology:virus}{vírus} não consegue mudar (\omterm{def:b3:adaptive-immunity:antibody}{anticorpos}
amplamente neutralizantes), depois de anos de maturação. O \omterm{def:b3:adaptive-immunity:germinal-centre}{centro
germinativo} pode ser guiado: uma sequência de imunógenos que primeiro
engaje os precursores germinais certos e depois recompense a ligação ao
sítio conservado poderia conduzir a maturação por esse caminho em meses,
e não em anos.
\textbf{19.} $p_{c} = 1 - 1/15 = 93.3\,\%$. Uma dose: $0.933/0.93 =
100.4\,\%$ --- inatingível; duas doses: $0.933/0.97 = 96.2\,\%$.
\textbf{20.} Imunes, $0.90\times 0.97 = 0.873$; $R_{\text{eff}} = 15\times
0.127 = 1.9$: o sarampo circula.
\textbf{21.} $60\,000\times (0.10 + 0.90\times 0.03) = 7600$ suscetíveis
por coorte de nascimento; depois de cinco anos, $38\,000$ --- uma
reserva na qual um caso importado encontra suscetíveis suficientes entre
seus contatos para sustentar cadeias.
\textbf{22.} Entre os suscetíveis, $s_{\infty} - \ln s_{\infty}/1.9 =
1$ dá $s_{\infty} \approx 0.23$: cerca de \qty{77}{\%} deles, uns
$29\,000$ casos.
\textbf{23.} Imunes, $0.95\times 0.97 = 0.92$; $R_{\text{eff}} = 15\times
0.08 = 1.2$ --- ainda acima de um; o limiar exige \qty{96}{\%} com duas
doses, e o sarampo é a doença que pune os últimos poucos por cento.
\textbf{24.} Os lactentes só são protegidos pela imunidade de quem está
em volta (e pelo \omterm{def:b3:adaptive-immunity:antibody}{anticorpo} materno por alguns meses): se as cadeias de
transmissão não conseguem se formar, nenhum caso chega até eles. Com
cobertura de \qty{85}{\%} a fração imune é $0.82$ e $R_{\text{eff}} = 2.6$: os surtos correm, e os lactentes, não vacinados e mais
vulneráveis, são infectados.
\textbf{25.} Repertório de cerca de $6\times 10^{10}$; \qty{2.7}{mg} de
IgG por dia dos \omterm{def:b3:adaptive-immunity:response}{plasmócitos} de uma resposta; cobertura de \qty{96}{\%}
com duas doses.
\end{solution}
